\section{Convergence and ergodicity}
\label{sec:convergence}

\subsection{Reduction to nonnegative centered moments}

\subsection{Late interactions}

\subsection{Relaxation of end patches}

\subsection{Spatial confinement and finite propagation}

\subsection{Common invariant limit}

\subsection{Uniform exponential ergodicity below the threshold}

\subsection{Uniform ergodicity under orientability}
\label{subsec:oriented-ergodicity}

The statement of Corollary~\ref{cor:oriented-ergodicity} is that orientability suffices to strengthen 
Theorem~\ref{thm:common-invariant-limit} to uniform ergodicity. The main idea is as follows: since all measures in~$\mathcal{M}_*$ converge to~$\pi$ then ,to prove ergodicity, it suffices to show that for any configuration~$\eta$ its evolution~$\delta_\eta P_t$ converges towards~$\mathcal{M}_*$. More precisely, by compactness and the Feller property, it
suffices to show that~$\delta_\eta P_t$ approaches~$\mathcal M_*$ on finite
sets, in the sense that
\begin{equation*}
\liminf_{t\to\infty}\inf_\eta P_t(\chi_A^*)(\eta)\geq0
\qquad\text{for every }A\Subset\Lambda.
\end{equation*}
To obtain this bound, choose a finite set~$I=\bar A$ as in the definition
of orientability and consider the dynamics of the spin system on~$I$, conditional on a particular realization of the trajectory on the (autonomous) set~$U(A)$. This is the original spin system, but now on a finite set and with time-inhomogenous boundary conditions. This restricted system is still patch positive and, with the original
centering profile~$p^*$, preserves the class~$M_*^I$ of probability measures on~$\cb{0,1}^I$ with non-negative centered moments. The boundary conditions are not homogenous in time, so the restricted system need not converge to an invariant measure, but it does converge towards the class~$\mathcal{M}_*^I$. Finiteness of~$I$ and presence of noise are crucial here, because they guarantee that a positive part of the measure enters~$\mathcal{M}_*^I$ after running the restricted process for any positive time. This proves the convergence towards~$\mathcal{M}_*$, in the above sense, and hence ergodicity.

\begin{proof}[Proof of Corollary~\ref{cor:oriented-ergodicity}]
    

Fix a nonempty finite set~$A$ and take~$I=\bar A$ and~$U=U(A)$ from
orientability. For~$\xi\in\cb{0,1}^U$, let~$\eta^{U,\xi}$ agree with~$\xi$
on~$U$ and with~$\eta$ elsewhere. Since~$I\cup U$ is autonomous, fixing
the spins on~$U$ defines a spin system on~$\cb{0,1}^I$, with generator
and rates
\begin{align*}
\Lcal^{I,\xi}f(\eta)&=\Lcal f(\eta^{U,\xi}),\\
\lambda_i^{I,\xi}(\eta)&=\lambda_i(\eta^{U,\xi}),
\qquad
\rho_i^{I,\xi}(\eta)=\rho_i(\eta^{U,\xi})
\quad(i\in I).
\end{align*}
Denote the restricted semigroup by~$P_t^{I,\xi}$.

With the original centering profile~$(p_i^\star)_{i\in I}$, define
\begin{equation*}
\mathcal M_*^I
=\cb{\mu\in\mathcal P\bb{\cb{0,1}^I}:
       \mu(\chi_B^*)\geq0\text{ for every }B\subseteq I}.
\end{equation*}
Lemma~\ref{lem:boundary-moment-preservation} in
Appendix~\ref{app:boundary-moment-preservation} gives
\begin{equation*}
\mathcal M_*^I P_t^{I,\xi}\subseteq\mathcal M_*^I.
\end{equation*}
This holds for every boundary condition~$\xi$. The restricted system is also patch positive, though we do not use this fact directly.

We next show that, for every initial configuration~$\eta$ and boundary conditions, after a unit of time a fixed positive part of the evolution law enters~$\mathcal{M}_*^I$. To be more precise we obtain a decomposition
\begin{equation}\label{eq:mass_leak_decomposition}
    \delta_\eta P_{s_1}^{I,\xi_1}\cdots P_{s_n}^{I,\xi_n} = \alpha_I \delta_{\zeta_I} + \bb{1-\alpha_I}\nu,
\end{equation}
where~$\zeta_I(i) = \ind\bb{p_i^* > 0}$ for~$i\in I$,~$\nu$ is a probability measure,~$\alpha_I>0$ and~$\xi_1,\dots,\xi_n$ and~$s_1+\dots +s_n\geq 1$ are arbitrary boundary conditions and their durations. This is useful because~$\delta_{\zeta_I}\in\mathcal{M}_*^I$.

We verify that the restricted process has positive probability of entering~$\zeta_I$ after any positive time. The rate at which a site turns to state~$1$ is~$\lambda_i\rho_i  \geq \varepsilon p_i^*$, which is positive if~$p_i^*>0$. Similarly, the rate at which a site turns to state~$0$ is~$\lambda_i\bb{1 - \rho_i}\geq \varepsilon \bb{1-p_i^*}$, which is positive if~$p_i^* = 0$. Both inequalities follow from patch positivity and uniform deaths. Thus every site~$i\in I$ has uniformly positive rate at which it takes the state~$\zeta_I(i)$. The event that every site in~$I$ updates during a unit interval and its last update sets it to~$\zeta_I$ therefore has probability bounded below by some~$\alpha_I\in(0,1)$, uniformly in the initial configuration and boundary trajectory. Applying this to the last unit interval gives~\eqref{eq:mass_leak_decomposition}.

Now we show that the process converges towards~$\mathcal{M}_*$ from any initial configuration~$\eta$. Define the~$\sigma$-algebra 
\begin{equation*}
\mathcal G_U:=\sigma\cb{\eta_s(j):j\in U,\ s\geq0}.
\end{equation*}
Since both~$U$ and~$I\cup U$ are autonomous,
conditional on~$\mathcal G_U$, the process on~$I$ evolves with
the revealed trajectory on~$U$ as a boundary condition. Only the states at~$\partial I:=N_*(I)\setminus I\subseteq U$ matter.
This set is finite and the update rates are bounded, so the boundary conditions change finitely many times before time~$t$. 
If the successive boundary configurations are~$\xi_1,\ldots,\xi_n$,
extended arbitrarily from~$\partial I$ to~$U$, and their durations
are~$t_1,\ldots,t_n$ with~$\sum_{k=1}^n t_k=t$, then, for~$f$ supported on~$I$,
\begin{equation*}
\CE{f(\eta_t)}{\mathcal G_U}
=P_{t_1}^{I,\xi_1}\cdots P_{t_n}^{I,\xi_n}(f)(\eta).
\end{equation*}
For~$t\geq1$, put~$m=\lfloor t\rfloor$ and split this evolution into~$m$
blocks: the first~$m-1$ have length~$1$, and the last has length~$t-m+1$.
Splitting the boundary intervals at the block endpoints gives
\begin{equation*}
\CE{f(\eta_t)}{\mathcal G_U}
=\bb{\prod_{j=1}^{m}
 \bb{P_{s_1^j}^{I,\xi_1^j}\cdots
     P_{s_{n_j}^j}^{I,\xi_{n_j}^j}}}(f)(\eta),
\end{equation*}
where~$\sum_l s_l^j=1$ for~$j<m$ and~$\sum_l s_l^m=t-m+1\geq1$;
products are ordered by increasing block index.

By~\eqref{eq:mass_leak_decomposition}, each block allows us to extract
mass~$\alpha_I$ at~$\zeta_I$ from the remaining measure.
Repeating the extraction gives
\begin{align*}
\CE{\chi_A^*(\eta_t)}{\mathcal G_U}
&=\sum_{k=1}^{m} \alpha_I(1-\alpha_I)^{k-1}
 \bb{\prod_{j=k+1}^{m}
  \bb{P_{s_1^j}^{I,\xi_1^j}\cdots
      P_{s_{n_j}^j}^{I,\xi_{n_j}^j}}}(\chi_A^*)(\zeta_I)\\
&\quad +(1-\alpha_I)^m\nu_t(\chi_A^*)\\
&\geq-\norm{\chi_A^*}_\infty(1-\alpha_I)^m,
\end{align*}
where~$\nu_t$ is a probability measure and the empty product is the identity.
The term indexed by~$k$ is the mass first extracted at the end of
block~$k$, evolved through the remaining blocks.
The sum is non-negative because~$\delta_{\zeta_I}\in\mathcal M_*^I$
and each~$P_s^{I,\xi}$ preserves~$\mathcal M_*^I$.
The lower bound side tends to zero exponentially,
and hence, for every~$A\Subset\Lambda$,
\begin{equation}\label{eq:convergence_towards_m*}
\liminf_{t\to\infty}\inf_\eta P_t(\chi_A^*)(\eta)\geq0.
\end{equation}

To obtain convergence to~$\pi$, we combine~\eqref{eq:convergence_towards_m*}
with Theorem~\ref{thm:common-invariant-limit}.
Fix a local function~$f$ and~$s>0$. Choose~$t_n\to\infty$ and
configurations~$\eta_n$ such that
$\abs{P_{t_n}f(\eta_n)-\pi(f)}$ converges to
$\limsup_{t\to\infty}\sup_\eta\abs{P_tf(\eta)-\pi(f)}$ and
$\delta_{\eta_n}P_{t_n-s}\Rightarrow\mu$ for some probability measure~$\mu$. This is possible because~$\mathcal{P}\bb{\cb{0,1}^\Lambda}$ is compact.
For every finite~$B\subseteq\Lambda$, the~\eqref{eq:convergence_towards_m*} implies
\begin{equation*}
\mu(\chi_B^*)
=\lim_{n\to\infty}P_{t_n-s}(\chi_B^*)(\eta_n)\geq0,
\end{equation*}
so~$\mu\in\mathcal M_*$. By the Feller property,~$P_sf$ is continuous,
and hence
\begin{equation*}
P_{t_n}f(\eta_n)
=(\delta_{\eta_n}P_{t_n-s})(P_sf)
\longrightarrow(\mu P_s)(f).
\end{equation*}
Theorem~\ref{thm:common-invariant-limit} now gives
\begin{equation*}
\limsup_{t\to\infty}\sup_\eta\abs{P_tf(\eta)-\pi(f)}
=\abs{(\mu P_s)(f)-\pi(f)}
\leq K_f(1+s)^D e^{-\varepsilon s/2}.
\end{equation*}
Letting~$s\to\infty$ proves convergence uniformly over initial
configurations for every local~$f$ and thus uniform ergodicity.

\par
\end{proof}
